% ============================================================================
% main.tex -- Documento principale (Progetto DashMe)
% Template report accademico con copertina dedicata, Abstract e capitoli
% ============================================================================
\documentclass[11pt,a4paper]{report}

% ============================================================================
% preamble.tex -- Pacchetti, palette cromatica, layout e macro personalizzate
% Progetto DashMe -- Università Politecnica delle Marche (UnivPM)
% Compatibile con Tectonic (XeTeX + BibTeX), pdfLaTeX, XeLaTeX e LuaLaTeX
% ============================================================================

\usepackage{iftex}
\ifPDFTeX
  \usepackage[utf8]{inputenc}
  \usepackage[T1]{fontenc}
  \usepackage{lmodern}
\else
  \usepackage{fontspec}
  \usepackage{lmodern}
\fi

\usepackage[italian]{babel}
\usepackage[a4paper, margin=2.3cm, headheight=14pt]{geometry}

\usepackage{graphicx}
\graphicspath{{../Tableau/}{figures/}}

\usepackage{booktabs}
\usepackage{tabularx}
\usepackage{array}
\usepackage{multirow}
\usepackage{amsmath}
\usepackage{amssymb}
\usepackage{xcolor}
\usepackage{caption}
\usepackage{subcaption}
\usepackage{float}
\usepackage{enumitem}
\usepackage{fancyhdr}
\usepackage{titlesec}
\usepackage{microtype}
\usepackage[numbers,sort&compress]{natbib}

% ----------------------------------------------------------------------------
% Palette cromatica accademica
% ----------------------------------------------------------------------------
\definecolor{primary}{HTML}{1B365D}
\definecolor{accent}{HTML}{D95F02}
\definecolor{muted}{HTML}{5A6577}
\definecolor{softbg}{HTML}{F4F6F9}
\definecolor{todobg}{HTML}{FFF3CD}
\definecolor{todoborder}{HTML}{B78103}

% ----------------------------------------------------------------------------
% Tipografia di titoli, didascalie ed elenchi
% ----------------------------------------------------------------------------
\titleformat{\section}{\large\bfseries\color{primary}}{\thesection}{0.75em}{}
\titleformat{\subsection}{\normalsize\bfseries\color{primary}}{\thesubsection}{0.6em}{}
\titleformat{\subsubsection}{\normalsize\bfseries\color{muted}}{\thesubsubsection}{0.6em}{}

\captionsetup{
  font=small,
  labelfont={bf,color=primary},
  labelsep=period
}

\setlist{nosep, leftmargin=1.5em}

% ----------------------------------------------------------------------------
% Intestazione e piè di pagina (fancyhdr)
% ----------------------------------------------------------------------------
\pagestyle{fancy}
\fancyhf{}
\fancyhead[L]{\footnotesize\color{muted}Analisi BI dei Richiami FDA -- Progetto DashMe}
\fancyhead[R]{\footnotesize\color{muted}UnivPM -- Project Management \& Business Intelligence}
\fancyfoot[C]{\footnotesize\color{muted}\thepage}
\renewcommand{\headrulewidth}{0.4pt}
\renewcommand{\headrule}{\hbox to\headwidth{\color{muted!40}\leaders\hrule height \headrulewidth\hfill}}
\renewcommand{\footrulewidth}{0pt}

\fancypagestyle{plain}{%
  \fancyhf{}%
  \fancyfoot[C]{\footnotesize\color{muted}\thepage}%
  \renewcommand{\headrulewidth}{0pt}%
  \renewcommand{\footrulewidth}{0pt}%
}

% ----------------------------------------------------------------------------
% Macro pubbliche personalizzate
% ----------------------------------------------------------------------------

% 1) \todo{testo}: pill/box ambra per annotazioni o completamenti pendenti
\newsavebox{\todobox}
\newcommand{\todo}[1]{%
  \sbox{\todobox}{\small\textbf{[DA COMPLETARE:} #1\textbf{]}}%
  \ifdim\wd\todobox>0.88\linewidth
    \par\smallskip\noindent
    {%
      \setlength{\fboxsep}{5pt}%
      \fcolorbox{todoborder}{todobg}{%
        \parbox{\dimexpr\linewidth-2\fboxsep-2\fboxrule\relax}{%
          \small\color{black}\textcolor{todoborder}{\textbf{[DA COMPLETARE:}} #1\textcolor{todoborder}{\textbf{]}}%
        }%
      }%
    }%
    \par\smallskip
  \else
    \fcolorbox{todoborder}{todobg}{%
      \small\color{black}\textcolor{todoborder}{\textbf{[DA COMPLETARE:}} #1\textcolor{todoborder}{\textbf{]}}%
    }%
  \fi
}

% 2) \dashboardfig[width]{filename}{caption}{label}:
%    cerca in ../Tableau/ e poi in figures/; se assente mostra un segnaposto pulito
\newcommand{\dashboardfig}[4][\textwidth]{%
  \begin{figure}[htbp]
    \centering
    \IfFileExists{../Tableau/#2}{%
      \includegraphics[width=#1,keepaspectratio]{../Tableau/#2}%
    }{%
      \IfFileExists{figures/#2}{%
        \includegraphics[width=#1,keepaspectratio]{figures/#2}%
      }{%
        {%
          \setlength{\fboxsep}{10pt}%
          \setlength{\fboxrule}{0.8pt}%
          \fcolorbox{muted}{softbg}{%
            \parbox{\dimexpr0.92\linewidth-2\fboxsep-2\fboxrule\relax}{%
              \centering\vspace{1.2em}
              {\large\bfseries\color{primary}[Figura segnaposto: \texttt{\detokenize{#2}}]}\\[0.6em]
              {\small\color{muted}Immagine non ancora esportata. Esportare la dashboard da Tableau e salvare il file in \texttt{docs/Tableau/\detokenize{#2}} (oppure in \texttt{docs/thesis/figures/\detokenize{#2}}), quindi ricompilare il documento.}%
              \vspace{1.2em}
            }%
          }%
        }%
      }%
    }%
    \caption{#3}%
    \label{#4}%
  \end{figure}%
}

% 3) \keyinsight{titolo}{corpo}: box di sintesi con sfondo softbg e bordo primary
\newcommand{\keyinsight}[2]{%
  \par\medskip\noindent
  {%
    \setlength{\fboxsep}{8pt}%
    \setlength{\fboxrule}{0.8pt}%
    \fcolorbox{primary}{softbg}{%
      \parbox{\dimexpr\linewidth-2\fboxsep-2\fboxrule\relax}{%
        {\color{primary}\bfseries #1}\par\smallskip
        {\small #2}%
      }%
    }%
  }%
  \par\medskip
}

% ----------------------------------------------------------------------------
% Hyperref (caricato per ultimo)
% ----------------------------------------------------------------------------
\PassOptionsToPackage{hyphens}{url}
\usepackage{hyperref}
\hypersetup{
  colorlinks=true,
  linkcolor=primary,
  citecolor=primary,
  urlcolor=primary,
  pdfborder={0 0 0},
  pdftitle={Analisi Business Intelligence dei Richiami FDA: Geografia Industriale, Profili di Rischio e Governance Operativa},
  pdfauthor={Gabriele Vianello},
  pdfsubject={Project Management e Business Intelligence -- Università Politecnica delle Marche},
  pdfkeywords={Business Intelligence, Visual Analytics, Tableau, Richiami FDA, Gestione del Rischio}
}


\begin{document}

\pagenumbering{Alph}
\begin{titlepage}
  \thispagestyle{empty}
  \centering

  \vspace*{0.4cm}
  {\Large\bfseries UNIVERSITÀ POLITECNICA DELLE MARCHE\par}
  \vspace{0.35cm}
  {\large Corso di Project Management \& Business Intelligence\par}

  \vfill

  {\huge\bfseries Analisi Business Intelligence dei Richiami FDA\par}
  \vspace{0.45cm}
  {\Large Geografia Industriale, Profili di Rischio e Governance Operativa\par}
  \vspace{0.55cm}
  {\normalsize\textit{Progetto DashMe}\par}

  \vfill

  \noindent
  \begin{minipage}[t]{0.46\textwidth}
    \raggedright
    {\bfseries Docenti del corso:}\\[0.4em]
    Prof.\ D.\ Ursino\\
    Prof.\ F.\ Parlapiano
  \end{minipage}%
  \hfill
  \begin{minipage}[t]{0.46\textwidth}
    \raggedleft
    {\bfseries Studente:}\\[0.4em]
    Gabriele Vianello
    % Aggiungere qui eventuali coautori e matricole
  \end{minipage}

  \vfill

  {\normalsize \today\par}
  \vspace*{0.2cm}
\end{titlepage}

\pagenumbering{roman}
\begin{abstract}
\noindent
\textbf{Contesto.} Il monitoraggio dei richiami di prodotto notificati dalla \emph{U.S. Food and Drug Administration} (FDA) rappresenta una sfida critica per la sicurezza dei consumatori, la continuità delle filiere produttive e la conformità regolatoria nei settori alimentare, farmaceutico, cosmetico e dei dispositivi medici.

\smallskip\noindent
\textbf{Obiettivo.} Il progetto \emph{DashMe} propone un sistema di Business Intelligence e Visual Analytics progettato per trasformare i dati grezzi degli eventi di richiamo FDA in viste decisionali integrate, capaci di supportare l'analisi geografica, l'individuazione dei profili di rischio clinico e il presidio operativo dei richiami ancora aperti sul mercato.

\smallskip\noindent
\textbf{Metodologia.} A partire dal dataset originale (\texttt{recalls.xlsx}) normalizzato in ambiente Tableau (\texttt{dati\_trasformati\_tableau.hyper}), sono state sviluppate regole di \emph{data cleansing} e campi calcolati in lingua italiana per la tassonomia settoriale, la classificazione delle cause di richiamo, il dettaglio degli allergeni, il packaging, le modalità di conservazione lungo la catena del freddo e la portata distributiva, articolando i risultati in tre dashboard analitiche.

\smallskip\noindent
\textbf{Principali Risultati.} Su un totale di 104.537 prodotti richiamati, l'analisi evidenzia una forte concentrazione territoriale in specifici hub statunitensi, profili di severità clinica marcatamente differenziati per settore merceologico (con incidenze elevate di Classe~I nei comparti alimentare e veterinario) e 18.555 richiami tuttora in corso che richiedono azioni mirate di governance operativa, controllo della catena del freddo e gestione del packaging.

\medskip\noindent
\textbf{Parole chiave:} Business Intelligence; Visual Analytics; Tableau; Richiami FDA; Gestione del Rischio.
\end{abstract}

% Decommentare la riga seguente per generare l'indice generale:
% \tableofcontents\clearpage

\pagenumbering{arabic}
\chapter{Introduzione}
\label{chap:introduction}

La sicurezza alimentare e la vigilanza sui prodotti regolamentati costituiscono un ambito critico di interesse pubblico e industriale, in cui la tutela della salute dei consumatori si intreccia strettamente con le discipline del \textit{Project Management} e del \textit{Supply Chain Risk Management}. Lungo filiere agroalimentari, farmaceutiche e manifatturiere sempre più interconnesse su scala transnazionale, un'anomalia microbiologica, un'adulterazione chimica o un difetto di etichettatura possono propagarsi rapidamente dai siti di produzione e dai varchi doganali fino alla distribuzione al dettaglio, generando impatti sanitari, costi logistici di ritiro, sanzioni regolatorie e danni reputazionali durevoli.

In questo scenario, le autorità di regolazione producono e pubblicano archivi pubblici di straordinaria ricchezza informativa, ma strutturalmente eterogenei per granularità, modello di rilevazione e finalità di intervento. Dalla prospettiva del governo del rischio, l'osservazione isolata di un singolo registro fornisce una visione parziale della vulnerabilità delle filiere. Al contrario, la lettura congiunta di più fonti istituzionali consente di presidiare l'intero ciclo di vita del rischio regolatorio: dalla sorveglianza chimica preventiva sui campioni di laboratorio, al blocco delle partite non conformi alle frontiere, fino alla gestione post-vendita dei richiami sul mercato interno.

Le metodologie di \textit{Business Intelligence} (BI) e \textit{Visual Analytics} offrono il supporto tecnico e metodologico necessario per trasformare questi volumi di dati grezzi e frammentati in strumenti decisionali interattivi, governati da pipeline riproducibili di \textit{Extract, Transform, Load} (ETL), modellazione dimensionale e progettazione visiva orientata alla chiarezza quantitativa \citep{sharda2020,few2013}. Su queste premesse nasce il progetto \textbf{DashMe}, sviluppato nell'ambito del corso di \textit{Project Management \& Business Intelligence} presso l'Università Politecnica delle Marche (UnivPM), che integra tre moduli analitici distinti realizzati con le tre principali piattaforme di BI contemporanee, seguendo la sequenza logica e architettonica dal livello associativo esplorativo fino alla modellazione relazionale avanzata:

\begin{enumerate}
    \item \textbf{Qlik Sense -- Sorveglianza e Allerta Rapida Europea (RASFF)}: il primo pilastro analizza le notifiche del sistema europeo \textit{Rapid Alert System for Food and Feed} \citep{rasff_portal}. Sfruttando il motore associativo in-memory di Qlik Sense, il modulo esamina \textbf{11.787 notifiche validate} (2012--2025) provenienti da 30 Paesi notificanti verso 141 Paesi di origine, concentrandosi su urgenza operativa (43,73\% di allarmi immediati), gravità clinica (66,8\% di rischio grave), respingimenti alla frontiera (10,6\%, pari a 1.255 casi) e prevalenza biologica (\emph{Salmonella spp.} prima causa con 2.893 segnalazioni).

    \item \textbf{Tableau -- Richiami Regolatori Statunitensi (FDA)}: il secondo pilastro adotta l'ambiente di \textit{Visual Analytics} di Tableau per esaminare i rapporti di \textit{enforcement} della \textit{U.S. Food and Drug Administration} \citep{fda_recalls_portal}. A partire da \textbf{104.537 prodotti richiamati} associati a \textbf{37.848 eventi univoci} nell'arco 2012--2026 e arricchiti mediante 14 campi calcolati in lingua italiana, il modulo indaga la geografia industriale statunitense, i profili di severità clinica nelle tre classi regolatorie FDA (Classe~I, Classe~II e Classe~III) \citep{fda_recall_classes}, l'epidemiologia di cause e allergeni nonché la governance operativa dei 18.555 richiami tuttora \emph{In corso}, della catena del freddo e del confezionamento (ove il packaging in vetro registra un'incidenza del 74,4\% di adulterazioni farmacologiche).

    \item \textbf{Microsoft Power BI -- Sorveglianza Chimica Preventiva (EFSA / WHO GEMS)}: il terzo pilastro affronta la prevenzione analitica a monte tramite i dati europei dell'archivio congiunto dell'Autorità Europea per la Sicurezza Alimentare e dell'Organizzazione Mondiale della Sanità \citep{who_gems_efsa}. Basato su uno \textit{Star Schema} composto da una tabella dei fatti e quattro tabelle dimensionali su \textbf{2.424.475 rilevazioni laboratoristiche} (1984--2017), il modulo impiega misure DAX rigorose per determinare il tasso di positività effettivo al netto dei 3.330 valori nulli (\textbf{39,63\%} globale; 43,44\% sui 32 Stati nazionali nella prima dashboard; 36,96\% sugli analiti specifici nella seconda), evidenziando criticità di bioaccumulo (65.604 positivi al cadmio nella carne e nelle frattaglie) e identificando i driver statistici della contaminazione mediante algoritmi di intelligenza artificiale visiva (\emph{Key Influencers}, fino a 2,53x).
\end{enumerate}

La Tabella~\ref{tab:intro_modules_overview} sintetizza le caratteristiche strutturali, i volumi informativi e l'articolazione in viste interattive dei tre moduli che compongono l'ecosistema DashMe.

\begin{table}[htbp]
  \centering
  \caption{Quadro sinottico dei tre moduli analitici del progetto DashMe.}
  \label{tab:intro_modules_overview}
  \small
  \begin{tabularx}{\textwidth}{@{}>{\raggedright\arraybackslash}p{2.2cm}>{\raggedright\arraybackslash}X>{\raggedright\arraybackslash}Xp{1.7cm}p{1.9cm}@{}}
    \toprule
    \textbf{Piattaforma BI} & \textbf{Fonte Dati / Ente} & \textbf{Granularità e Volume} & \textbf{Arco Temp.} & \textbf{N. Dashboard} \\
    \midrule
    Qlik Sense &
    Portale RASFF / Commissione Europea (\texttt{RASFF\_\-window.csv}) &
    11.787 notifiche validate (30 Paesi notificanti, 141 Paesi di origine) &
    2012--2025 &
    4 fogli \\
    \addlinespace[0.4em]
    Tableau &
    FDA Enforcement Reports (\texttt{recalls.xlsx}) &
    104.537 prodotti in 37.848 eventi (14 campi calcolati) &
    2012--2026 &
    3 dashboard \\
    \addlinespace[0.4em]
    Power BI &
    EFSA / WHO GEMS (\texttt{Contamination\_\-european\_\-food.csv}) &
    2.424.475 rilevazioni (33 entità geografiche, 50 contaminanti, 348 alimenti; Star Schema 1 Fact + 4 Dim) &
    1984--2017 &
    3 dashboard \\
    \bottomrule
  \end{tabularx}
\end{table}

Alla luce di tale architettura multi-piattaforma, il presente elaborato si propone di rispondere a quattro domande di ricerca (\textit{Research Questions}, RQ):

\begin{enumerate}
    \item \textbf{RQ1 (Qlik Sense -- Sorveglianza e allerta rapida UE RASFF)}: quali trend storici, flussi geopolitici di confine, matrici alimentare-pericolo e dinamiche di tracciabilità emergono dalle notifiche RASFF?
    \item \textbf{RQ2 (Tableau -- Richiami regolatori USA FDA)}: come si articolano geografia industriale, profili di severità clinica (Classi I--III), epidemiologia delle cause/allergeni e governance operativa (catena del freddo, packaging, portata distributiva) nei richiami FDA?
    \item \textbf{RQ3 (Power BI -- Sorveglianza chimica europea EFSA / WHO GEMS)}: quali relazioni intercorrono tra volume di campionamento, tassi di positività al netto dei valori nulli, fasce di concentrazione per filiera, screening multi-analita e fattori di influenza AI?
    \item \textbf{RQ4 (Valutazione comparativa multi-piattaforma)}: in che misura Qlik Sense (motore associativo), Tableau (visual analytics esplorativa) e Power BI (modellazione stellare VertiPaq, DAX e AI visiva) offrono vantaggi e compromessi complementari?
\end{enumerate}

\section*{Struttura dell'Elaborato}
Il resto della trattazione è articolato in sei capitoli successivi:
\begin{itemize}
    \item il \textbf{Capitolo~\ref{chap:methodology}} illustra l'architettura metodologica complessiva del progetto DashMe, descrivendo le basi dati sorgente, le pipeline di \emph{data cleansing} e \emph{feature engineering} e i criteri di progettazione visiva adottati sulle tre piattaforme;
    \item il \textbf{Capitolo~\ref{chap:qlik}} presenta i risultati dell'analisi in \emph{Qlik Sense} sul sistema europeo di allerta rapida RASFF lungo quattro fogli dedicati a sintesi esecutiva, geopolitica dei controlli, matrici alimentari e tracciabilità di filiera;
    \item il \textbf{Capitolo~\ref{chap:tableau}} approfondisce le evidenze prodotte in \emph{Tableau} sui richiami regolatori statunitensi FDA, coprendo geografia industriale e responsabilità aziendale, epidemiologia delle cause con focus sugli allergeni e governance operativa su logistica e packaging;
    \item il \textbf{Capitolo~\ref{chap:powerbi}} esamina il modulo \emph{Microsoft Power BI} per la sorveglianza chimica europea EFSA / WHO GEMS, analizzando il paradigma volume-positività, il profilo di rischio per filiera e fasce di concentrazione, le dinamiche multi-analita e i fattori determinanti individuati tramite AI visiva;
    \item il \textbf{Capitolo~\ref{chap:discussion}} propone una discussione trasversale delle implicazioni manageriali per il governo delle filiere, confronta criticamente i tre strumenti di BI rispetto alle capacità tecniche osservate e discute i limiti metodologici dei dataset analizzati;
    \item il \textbf{Capitolo~\ref{chap:conclusions}} sintetizza le risposte alle quattro domande di ricerca e delinea le possibili estensioni future dell'ecosistema DashMe.
\end{itemize}

\chapter{Metodologia, Architetture Dati e Progettazione Visiva}
\label{chap:methodology}

L'architettura analitica del progetto DashMe integra tre piattaforme di Business Intelligence --- \textit{Qlik Sense}, \textit{Tableau} e \textit{Microsoft Power BI} --- applicate rispettivamente alla sorveglianza europea delle allerte alimentari (RASFF), al monitoraggio regolatorio statunitense dei richiami di prodotto (FDA) e alla rilevazione laboratoristica preventiva dei contaminanti chimici in Europa (EFSA / WHO GEMS). Ciascun modulo adotta il paradigma di modellazione ed estrazione più aderente alla struttura del proprio dataset sorgente, mantenendo al contempo criteri condivisi di pulizia del dato, normalizzazione linguistica in italiano e progettazione visiva \citep{kimball2013,few2013}.

% ============================================================================
\section{Pipeline Qlik Sense e Modello Associativo RASFF}
\label{sec:method_qlik}
% ============================================================================

Il primo flusso metodologico opera sul dataset europeo del sistema di allerta rapida per alimenti e mangimi RASFF (\texttt{data/raw/RASFF\_window.csv}) \citep{rasff_portal}. Il tracciato grezzo si compone di \textbf{11.788} record descritti da \textbf{14} colonne originali: \texttt{reference}, \texttt{category}, \texttt{type}, \texttt{subject}, \texttt{date}, \texttt{notifying\_country}, \texttt{classification}, \texttt{risk\_decision}, \texttt{distribution}, \texttt{forAttention}, \texttt{forFollowUp}, \texttt{operator}, \texttt{origin} e \texttt{hazards}. L'arco temporale complessivo si estende dal \textbf{24/01/2012} al \textbf{31/12/2025}, con oltre il \textbf{98,8\%} delle notifiche concentrate nel periodo \textbf{2020--2025}.

\subsection{Profilazione Qualitativa e Gestione delle Anomalie CSV}
L'ispezione preliminare del file sorgente ha evidenziato una specifica anomalia sintattica su \textbf{4} righe del tracciato (\texttt{2024.6416}, \texttt{2024.4612}, \texttt{2023.7176} e \texttt{2021.2621}): la presenza di doppi apici non correttamente gestiti all'interno del campo testuale \texttt{subject} provocava, in fase di lettura grezza, uno slittamento di colonna verso destra che trasferiva frammenti descrittivi dentro il campo \texttt{notifying\_country}, generando false entità nazionali. Il trattamento di parsing e validazione implementato nello script di caricamento Qlik Sense ha risolto tali disallineamenti consolidando un perimetro valido pari a \textbf{11.787} notifiche distinte (\textbf{11,79k}), \textbf{30} Paesi/entità notificanti effettive (i 27 Stati membri dell'Unione Europea affiancati dai Paesi EFTA e dalla Commissione Europea) e \textbf{141} Paesi di origine delle merci.

\subsection{Modello Associativo, Normalizzazione \texttt{1:N} e Set Analysis}
In coerenza con il motore associativo in-memory di Qlik Sense \citep{qlik_docs}, il modello evita la denormalizzazione piatta dei campi multivalore attraverso una struttura stellata centrata sulla tabella dei fatti \texttt{[Notifiche]}, identificata dalla chiave primaria testuale \texttt{NotificationID = Text(reference)}. I campi contenenti liste separate da delimitatori vengono scomposti mediante la funzione \texttt{SubField()} in tabelle ponte con cardinalità \texttt{1:N}:
\begin{itemize}
    \item \texttt{[Bridge\_Origin]}: disaggrega il campo \texttt{origin} per associare correttamente ciascuna notifica a uno o più Paesi di provenienza della partita;
    \item \texttt{[Bridge\_Hazards]}: espande il campo \texttt{hazards} separando i singoli pericoli segnalati ed estraendo in modo strutturato la \texttt{SostanzaContaminante} e la relativa \texttt{CategoriaPericolo} tramite espressioni \texttt{TextBetween(..., '\{', '\}')};
    \item \textbf{Tabella di distribuzione}: normalizza l'elenco degli Stati coinvolti nella catena distributiva lungo la dimensione \texttt{PaeseDestinazione}.
\end{itemize}
Questa scomposizione relazionale consente l'esplorazione associativa bidirezionale tra qualsiasi Paese di origine, contaminante e Stato destinatario, prevenendo al contempo il doppio conteggio volumetrico nei KPI grazie all'aggregazione sistematica su \texttt{Count(DISTINCT NotificationID)}.

Sul piano delle trasformazioni semantiche e geospaziali, lo script di caricamento applica quattro interventi principali:
\begin{enumerate}
    \item \textbf{Localizzazione dei domini}: traduzione integrale in italiano delle dimensioni \texttt{LivelloRischio}, \texttt{TipoNotifica}, \texttt{TipoProdotto}, \texttt{Categoria Alimento} e delle denominazioni nazionali;
    \item \textbf{Separazione geospaziale}: distinzione esplicita tra l'attributo descrittivo \texttt{PaeseOrigine} (che valorizza i casi non attribuiti come \texttt{"Origine sconosciuta"}) e il campo cartografico \texttt{PaeseOrigine\_Geometria} (impostato a \texttt{Null} in assenza di origine certa, al fine di impedire geocoding errati sulle mappe);
    \item \textbf{Consolidamento microbiologico}: unificazione delle molteplici varianti e sierotipi di Salmonella mediante la regola \texttt{WildMatch(Lower(\-SostanzaContaminante), '*salmonella*')} $\to$ \texttt{'Salmonella spp.'};
    \item \textbf{Espressioni di Set Analysis}: formalizzazione dinamica degli indicatori decisionali, tra cui \texttt{\% Allarmi Immediati}, \texttt{\% Rischio Grave}, il confronto tra partite \texttt{Bloccate in Dogana} e \texttt{Distribuite a Scaffale}, nonché le \texttt{Richieste Follow-Up}.
\end{enumerate}

% ============================================================================
\section{Pipeline Tableau e Feature Engineering sui Richiami FDA}
\label{sec:method_tableau}
% ============================================================================

Il secondo flusso metodologico esamina i richiami regolatori statunitensi pubblicati dalla \textit{U.S. Food and Drug Administration} (FDA) a partire dal file sorgente \texttt{data/recalls.xlsx} \citep{fda_recalls_portal,tableau_docs}. Il dataset comprende esattamente \textbf{104.537} righe --- ciascuna associata a un \texttt{Product ID} univoco --- descritte da \textbf{17} colonne originali: \texttt{FEI Number}, \texttt{Recalling Firm Name}, \texttt{Recalling Firm City}, \texttt{Recalling Firm State}, \texttt{Recalling Firm Country}, \texttt{Event ID}, \texttt{Event Classification}, \texttt{Center}, \texttt{Center Classification Date}, \texttt{Status}, \texttt{Recall Details}, \texttt{Product ID}, \texttt{Product Type}, \texttt{Product Classification}, \texttt{Product Description}, \texttt{Reason for Recall} e \texttt{Distribution Pattern}.

La profilazione relazionale sull'arco temporale \textbf{08/06/2012}--\textbf{09/10/2026} evidenzia \textbf{37.848} \texttt{Event ID} distinti (con una molteplicità media pari a \textbf{2,76} prodotti per evento), \textbf{10.276} identificativi di stabilimento \texttt{FEI Number}, \textbf{10.991} ragioni sociali grezze poi consolidate in \textbf{10.113} entità aziendali pulite, \textbf{33.701} descrizioni causali distinte e \textbf{22.485} pattern distributivi distinti. Sul versante geografico sono rappresentati \textbf{57} Paesi, con \textbf{99.130} record attribuiti agli Stati Uniti (\textbf{94,83\%}) e \textbf{5.407} record extra-USA (\textbf{5,17\%}); all'interno del perimetro statunitense, la normalizzazione produce \textbf{53} entità statali e territoriali pulite in \texttt{Stato Notificante (Pulito)}.

\subsection{Granularità della Severità e Architettura dei Workbook}
Un passaggio metodologico essenziale riguarda la distinzione tra la severità definita a livello di singolo articolo (\texttt{Product Classification}) e quella attribuita all'intero procedimento amministrativo (\texttt{Event Classification}). I due campi divergono su \textbf{2.735} righe del dataset; poiché la metrica volumetrica delle viste adotta come unità elementare il singolo prodotto richiamato (\texttt{COUNT([Product ID])}), l'analisi di letalità utilizza coerentemente \texttt{Product Classification} \citep{fda_recall_classes}, ripartita in \textbf{Classe~I = 20.165} (\textbf{19,29\%}), \textbf{Classe~II = 75.339} (\textbf{72,07\%}) e \textbf{Classe~III = 9.033} (\textbf{8,64\%}).

Sotto il profilo degli artefatti software, la cartella di lavoro \texttt{src/Tableau/dashsmash3.twb} conserva l'intero modello ETL composto dai \textbf{14 campi calcolati} (corredati da \textbf{53} alias di dominio) e i 4 fogli costitutivi della Dashboard~3, mentre le Dashboard~1 e~2 sono state realizzate su cartelle Tableau distinte condividendo esattamente il medesimo schema di campi calcolati.

\subsection{Tassonomia dei 14 Campi Calcolati e Regole Lessicali}
La Tabella~\ref{tab:calculated_fields} riepiloga i \textbf{14 campi calcolati} implementati in \texttt{src/Tableau/dashsmash3.twb} per trasformare i testi liberi FDA in dimensioni analitiche strutturate.

\begin{table}[htbp]
\centering
\caption{Riepilogo dei 14 campi calcolati implementati in Tableau per pulizia, normalizzazione e \textit{feature engineering} sui richiami FDA.}
\label{tab:calculated_fields}
\small
\setlength{\tabcolsep}{4.5pt}
\begin{tabularx}{\textwidth}{@{}>{\raggedright\arraybackslash}p{4.6cm}>{\raggedright\arraybackslash}p{2.7cm}>{\raggedright\arraybackslash}X@{}}
\toprule
\textbf{Campo Calcolato} & \textbf{Dimensione} & \textbf{Funzione nel Modello ETL e Analitico} \\
\midrule
\texttt{Conteggio prodotti} & Metrica volumetrica & Calcola \texttt{COUNT([Product ID])} per quantificare i singoli articoli richiamati. \\
\texttt{Ragione Sociale (Pulita)} & Anagrafica aziende & Consolida le 10.991 denominazioni grezze di \texttt{Recalling Firm Name} in 10.113 entità pulite. \\
\texttt{Stato Notificante (Pulito)} & Geografia & Standardizza \texttt{Recalling Firm State} in 53 entità statali e territoriali USA valide. \\
\texttt{Città Notificante (IT)} & Geografia & Normalizza e adatta per la vista italiana le località di \texttt{Recalling Firm City}. \\
\texttt{Tipo Prodotto (Atomico)} & Merceologia & Separa \texttt{Food/Cosmetics} (29.907) in \texttt{Cibo e Bevande} e \texttt{Cosmetici} con 11 parole chiave. \\
\texttt{Categoria Ragione di Richiamo} & Tassonomia causale & Classifica \texttt{Reason for Recall} tramite cascata \texttt{IF/ELSEIF} su 6 macro-cause più \texttt{Altro}. \\
\texttt{Dettaglio Allergene} & Sicurezza alimentare & Estrae la tipologia di allergene non dichiarato mediante flag lessicali dedicati. \\
\texttt{Numero Allergeni Coinvolti} & Sicurezza alimentare & Somma i flag binari degli allergeni rilevati per singolo prodotto richiamato. \\
\texttt{Tipo Packaging} & Confezionamento & Identifica 4 famiglie di materiale da \texttt{Product Description} tramite parole chiave. \\
\texttt{Modalità Conservazione} & Logistica e freddo & Classifica il regime termico (\texttt{Surgelato}, \texttt{Refrigerato / Fresco}, \texttt{Temperatura Ambiente / Altro}). \\
\texttt{Portata Distributiva} & Estensione mercato & Classifica l'ampiezza della distribuzione da \texttt{Distribution Pattern}. \\
\texttt{Distribuzione - Suddivisione 1, 2, 3} (3 campi) & Parsing logistico & Tre campi dedicati alla segmentazione progressiva di \texttt{Distribution Pattern} per aree servite. \\
\bottomrule
\end{tabularx}
\end{table}

L'audit delle espressioni condizionali sull'intero universo di \textbf{104.537} record consente di documentare con precisione le regole lessicali adottate e l'incidenza esatta delle classi residuali:
\begin{itemize}
    \item \textbf{\texttt{Tipo Prodotto (Atomico)}}: scomposizione della categoria composita originale \texttt{Food/Cosmetics} (\textbf{29.907} record) mediante \textbf{11 parole chiave cosmetiche}, che isolano esattamente \textbf{29.737} prodotti in \texttt{Cibo e Bevande} e \textbf{170} in \texttt{Cosmetici}.
    \item \textbf{\texttt{Categoria Ragione di Richiamo}}: cascata ordinata \texttt{IF/ELSEIF} applicata a \texttt{LOWER([Reason for Recall])}, nella quale le regole per l'individuazione di principi attivi farmaceutici sono valutate \emph{prima} del controllo generico sulla stringa \texttt{"undeclared"} per evitare falsi positivi nella classe degli allergeni. Le \textbf{6} macro-cause tracciate coprono \textbf{43.856} record (\textbf{41,95\%}), mentre la classe residuale \texttt{Altro} raccoglie \textbf{60.681} record (\textbf{58,05\%}).
    \item \textbf{\texttt{Tipo Packaging}}: ricerca di parole chiave su \texttt{LOWER([Product Description])} articolata in quattro classi materiali (\texttt{glass}/\texttt{jar}/\texttt{bottle}; \texttt{plastic}/\texttt{pouch}/\texttt{bag}; \texttt{can}/\texttt{aluminum}/\texttt{metal}; \texttt{box}/\texttt{carton}/\texttt{paper}). I \textbf{4} materiali identificati coprono \textbf{28.995} record (\textbf{27,74\%}), lasciando nella classe residuale \texttt{Altro / Non specificato} i restanti \textbf{75.542} record (\textbf{72,26\%}).
    \item \textbf{\texttt{Modalità Conservazione}}: ricerca lessicale su \texttt{LOWER([Product Description])} che mappa \texttt{frozen}/\texttt{ice cream} su \texttt{Surgelato} (\textbf{3.849} record, \textbf{3,68\%}) e \texttt{refrigerat}/\texttt{keep cold}/\texttt{fresh} su \texttt{Refrigerato / Fresco} (\textbf{3.884} record, \textbf{3,72\%}); la classe residuale \texttt{Temperatura Ambiente / Altro} assorbe \textbf{96.804} record (\textbf{92,60\%}).
    \item \textbf{\texttt{Dettaglio Allergene} e \texttt{Numero Allergeni Coinvolti}}: la classificazione combina la somma di \textbf{7 flag binari} (\texttt{milk|dairy}, \texttt{peanut}, \texttt{soy}, \texttt{wheat|gluten}, \texttt{egg}, \texttt{nut|almond|\-walnut|\-cashew|\-pecan}, \texttt{sesame}) con ulteriori controlli dedicati a pesce/crostacei e solfiti. Sull'intero dataset, la classe \texttt{Non applicabile / Altro} comprende \textbf{94.478} record (\textbf{90,38\%}). Sotto il profilo tecnico, va segnalata la sovrapposizione lessicale tra la sottostringa \texttt{"peanut"} e la sottostringa \texttt{"nut"} nel flag della frutta a guscio, che eleva sistematicamente a $\ge 2$ il punteggio calcolato in \texttt{Numero Allergeni Coinvolti} per ogni richiamo contenente arachidi.
\end{itemize}

% ============================================================================
\section{Pipeline Power BI, Star Schema e Modello DAX EFSA / WHO GEMS}
\label{sec:method_powerbi}
% ============================================================================

Il terzo modulo del progetto opera sul dataset di sorveglianza chimica europea EFSA / WHO GEMS (\texttt{data/raw/Contamination\_european\_food.csv}) \citep{who_gems_efsa}. L'archivio sorgente comprende \textbf{2.424.475} determinazioni laboratoristiche su \textbf{18} colonne originali riferite all'arco temporale \textbf{1984--2017}, con \textbf{33} entità geografiche, \textbf{50} contaminanti, \textbf{22} gruppi alimentari e \textbf{348} alimenti distinti. Il file CSV grezzo, pari a \textbf{516,6~MB}, viene compresso a \textbf{69,3~MB} all'interno del motore colonnare VertiPaq del file \texttt{src/PowerBI/powerdash.pbix}, conseguendo una riduzione di occupazione su disco dell'\textbf{86,6\%}.

\subsection{Star Schema a 1 Tabella dei Fatti e 4 Dimensioni in Power Query M}
Seguendo la metodologia dimensionale di Kimball \citep{kimball2013,powerbi_docs}, gli script Power Query M archiviati nella cartella \texttt{src/PowerBI/queries/} implementano uno \textbf{Star Schema} puro composto da una tabella dei fatti e quattro tabelle dimensionali collegate tramite relazioni \texttt{1:N}:
\begin{enumerate}
    \item \textbf{\texttt{Fact\_Contaminazioni}} (\texttt{Fact\_Contaminazioni.m}): \textbf{2.424.475} righe $\times$ \textbf{15} colonne. Lo script rimuove 6 colonne testuali interamente ridondanti o delegate alle dimensioni (\texttt{GEMSFoodName}, \texttt{RegionCode}, \texttt{UnitName}, \texttt{CountryName}, \texttt{FoodGroupName}, \texttt{ContaminantName}), applica la tipizzazione \texttt{Int64} agli identificativi e \texttt{type number} con locale \texttt{"en-US"} su \texttt{ResultValue}, normalizza \texttt{SerialNumber} (rimozione degli spazi superflui e conversione delle stringhe vuote in \texttt{null}) e genera tre colonne analitiche derivate: \texttt{Esito Analisi}, il flag binario \texttt{Is\_Positive} e la discretizzazione logaritmica \texttt{Fascia Concentrazione} in \textbf{7} classi (\texttt{0..6}, da \texttt{0. Dato Mancante} fino a \texttt{6. Critica (\textgreater{}1.000)}).
    \item \textbf{\texttt{Dim\_Geografia}} (\texttt{Dim\_Geografia.m}): \textbf{33} righe $\times$ \textbf{5} colonne. Separa esplicitamente i \textbf{32} Stati nazionali dall'aggregato sovranazionale \texttt{WHO European Region} (\texttt{CountryID = 333}) attraverso l'attributo \texttt{Tipo Entità Geografica} e imposta \texttt{Country\_Geo\_EN = null} per la voce OMS onde prevenire errate geocodificazioni cartografiche.
    \item \textbf{\texttt{Dim\_Contaminanti}} (\texttt{Dim\_Contaminanti.m}): \textbf{50} righe $\times$ \textbf{5} colonne. Traduce le denominazioni chimiche in italiano, aggrega gli analiti in \textbf{5} classi tramite \texttt{Famiglia Contaminante} e introduce il flag \texttt{Contaminante Specifico} per isolare la voce generica \texttt{Other}.
    \item \textbf{\texttt{Dim\_GruppiAlimentari}} (\texttt{Dim\_GruppiAlimentari.m}): \textbf{22} righe $\times$ \textbf{4} colonne. Affianca alla descrizione completa \texttt{Gruppo Alimentare (IT)} l'etichetta compatta \texttt{Gruppo Alimentare (breve)}, introdotta nell'ultimo affinamento progettuale per evitare il troncamento dei nodi nell'albero di scomposizione (\textit{Decomposition Tree}).
    \item \textbf{\texttt{Dim\_Alimenti}} (\texttt{Dim\_Alimenti.m}): \textbf{348} righe $\times$ \textbf{3} colonne (\texttt{FoodID}, \texttt{FoodName\_EN} e \texttt{Alimento (IT)}).
\end{enumerate}

\subsection{Gestione dei Valori Nulli in \texttt{ResultValue} e Modello DAX}
La verifica puntuale sul campo \texttt{ResultValue} evidenzia tre casistiche mutuamente esclusive sui \textbf{2.424.475} record del CSV: \textbf{959.424} determinazioni positive (\texttt{\textgreater{} 0}), \textbf{1.461.721} determinazioni negative sotto la soglia di rilevabilità (\texttt{= 0}, \texttt{\textless{} LOD}) e \textbf{3.330} record con risultato mancante (\texttt{null}, pari allo \textbf{0,14\%}). Poiché su un record privo di valore non è possibile stabilire l'esito positivo o negativo del test, la misura DAX \texttt{Tasso di Positività \%} esclude esplicitamente i \textbf{3.330} valori nulli dal denominatore mediante il filtro \texttt{NOT(ISBLANK(Fact\_Contaminazioni[ResultValue]))}:
\begin{equation}
\text{Tasso di Positività \%} = \frac{959.424}{2.421.145} = 39,63\%,
\end{equation}
evitando la sottostima che deriverebbe dall'inclusione delle righe nulle al denominatore ($959.424 / 2.424.475 = 39,57\%$).

Accanto a \texttt{Tasso di Positività \%}, il livello semantico DAX definisce le misure principali \texttt{Totale Analisi}, \texttt{Campioni Fisici Distinti}, \texttt{Analisi Positive (\textgreater{} LOD)}, \texttt{Concentrazione P95 Positivi ($\mu\text{g/kg}$)} e \texttt{Positività Benchmark Generale \%}, nonché le due colonne calcolate su \texttt{SerialNumber} --- \texttt{Analiti Testati sul Campione} e \texttt{Fascia Pannello Multi-Analita} --- necessarie per quantificare la densità di screening multi-residuo per singolo campione fisico.

% ============================================================================
\section{Confronto Architetturale e Principi di Visual Analytics}
\label{sec:method_comparison}
% ============================================================================

La Tabella~\ref{tab:methodology_comparison} sintetizza le scelte architetturali delle tre pipeline, evidenziando come il passaggio progressivo da dataset di medie dimensioni e schema semi-strutturato (RASFF) a registri amministrativi ad alto contenuto testuale (FDA) e archivi laboratoristici multi-milionari (EFSA / WHO GEMS) abbia guidato la scelta degli strumenti e dei modelli dati.

\begin{table}[htbp]
\centering
\caption{Confronto architetturale tra le tre pipeline Business Intelligence sviluppate nel progetto DashMe.}
\label{tab:methodology_comparison}
\small
\setlength{\tabcolsep}{4pt}
\begin{tabularx}{\textwidth}{@{}>{\raggedright\arraybackslash}p{2.7cm}>{\raggedright\arraybackslash}X>{\raggedright\arraybackslash}X>{\raggedright\arraybackslash}X@{}}
\toprule
\textbf{Criterio} & \textbf{Qlik Sense (RASFF)} & \textbf{Tableau (FDA)} & \textbf{Power BI (EFSA / OMS)} \\
\midrule
\textbf{Sorgente} & \texttt{data/raw/RASFF\_window.csv} (14 colonne) & \texttt{data/recalls.xlsx} (17 colonne) & \texttt{data/raw/Contamination\_} \texttt{european\_food.csv} (18 colonne) \\
\addlinespace[2pt]
\textbf{Volume dati} & 11.788 grezzi $\to$ 11.787 notifiche (2012--2025) & 104.537 prodotti / 37.848 eventi (2012--2026) & 2.424.475 analisi; CSV 516,6~MB $\to$ 69,3~MB VertiPaq ($-86,6\%$) \\
\addlinespace[2pt]
\textbf{Paradigma modello} & Associativo \texttt{1:N} con tabelle ponte (\texttt{[Bridge\_Origin]}, \texttt{[Bridge\_Hazards]}) & Tabella analitica arricchita (14 campi calcolati, 53 alias) & Star Schema (1 Fact + 4 Dimensioni) con misure DAX \\
\addlinespace[2pt]
\textbf{Pipeline ETL} & Script Qlik: \texttt{SubField()}, \texttt{TextBetween()}, \texttt{WildMatch()}, Set Analysis & Campi calcolati Tableau: \texttt{IF/ELSEIF} lessicali e flag binari & 5 script Power Query M (\texttt{Fact\_*} e \texttt{Dim\_*}) + colonne su \texttt{SerialNumber} \\
\addlinespace[2pt]
\textbf{Qualità dati} & Isolamento 4 righe CSV corrotte; \texttt{Null} geometrici per origine ignota & Pulizia 10.991 $\to$ 10.113 aziende; distinzione \texttt{Product} vs \texttt{Event Class.} & Esclusione 3.330 \texttt{null} dal denominatore ($39,63\%$); separazione aggregato OMS \\
\addlinespace[2pt]
\textbf{Numero di viste} & 4 fogli interattivi (Dashboard~1--4) & 3 dashboard analitiche (4 fogli nella Dashboard~3 in \texttt{dashsmash3.twb}) & 3 pagine/dashboard (incluse AI/Key Influencers e Decomposition Tree) \\
\bottomrule
\end{tabularx}
\end{table}

Nonostante la diversità tecnologica dei motori sottostanti, la progettazione dell'interfaccia visiva in tutti e tre i moduli segue principi unitari di \textit{Visual Analytics} \citep{few2013}:
\begin{itemize}
    \item \textbf{Separazione contestuale multi-dashboard}: anziché concentrare indicatori eterogenei in un'unica schermata sovraccarica, ciascun modulo articola l'esplorazione in tre o quattro viste specializzate dedicate progressivamente al quadro sintetico/geopolitico, al profilo di rischio per filiera o causa e alle dinamiche operative o statistiche di dettaglio.
    \item \textbf{Codifica cromatica coerente della gravità e dell'esito}: il colore non svolge una funzione decorativa ma veicola sistematicamente il gradiente di criticità clinica o analitica (livelli di rischio RASFF, Classe~I--III FDA, esiti \texttt{\textless{} LOD} vs \texttt{\textgreater{} LOD} e fasce di concentrazione EFSA), assicurando continuità cognitiva tra viste adiacenti.
    \item \textbf{Allineamento tra geometria del grafico e domanda analitica}: i confronti di ranking e incidenza settoriale privilegiano barre ordinate e barre impilate al 100\%, la distribuzione territoriale adotta mappe coropletiche o proporzionali depurate dalle entità prive di geometria certa (\texttt{Null}), le correlazioni tra volume e severità ricorrono a grafici a bolle/dispersione con soglie di riferimento, mentre l'indagine gerarchica e causale sfrutta rispettivamente alberi di scomposizione e componenti predittivi integrati.
\end{itemize}

\keyinsight{Sintesi Metodologica -- Complementarità delle Tre Architetture}{L'impiego combinato di Qlik Sense, Tableau e Power BI dimostra come la scelta dell'architettura dati dipenda direttamente dalla struttura del problema: il modello associativo con tabelle ponte \texttt{1:N} risolve nativamente la cardinalità multipla di origini e contaminanti RASFF senza doppio conteggio; i campi calcolati lessicali in Tableau trasformano oltre centomila descrizioni testuali FDA in tassonomie di rischio, allergeni e packaging; infine, lo Star Schema Power Query M e la compressione colonnare VertiPaq in Power BI consentono di interrogare oltre 2,42 milioni di analisi EFSA garantendo rigore statistico nella gestione dei valori nulli.}

\chapter{Analisi delle Dashboard del Progetto DashMe}
\label{sec:dashboard_analysis}

\section{Dashboard 1: Geografia Industriale, Settori Merceologici e Responsabilità Aziendale}
\label{subsec:dash1}

La prima dashboard (\textit{``Geografia Industriale, Settori Merceologici e Responsabilità Aziendale''}, Figura~\ref{fig:dash1}) risponde alla domanda \textbf{RQ1}, integrando la dimensione spaziale statunitense con il profilo di pericolosità dei settori merceologici e la concentrazione dei volumi presso i singoli operatori.

\dashboardfig{Dash1.jpeg}{Dashboard 1: Geografia Industriale, Settori Merceologici e Responsabilità Aziendale.}{fig:dash1}

\subsection{Concentrazione Territoriale negli Stati Uniti}
La vista cartografica ``Concentrazione Territoriale USA per Stato'' evidenzia una distribuzione fortemente asimmetrica degli eventi di richiamo sul territorio statunitense, con un intervallo che spazia da 1 fino a 4.522 eventi per singolo Stato. I principali poli di concentrazione emergono in California, Texas, Florida e lungo le direttrici industriali e logistiche del Nord-Est e del Midwest, cui si affianca una presenza rilevante a Puerto Rico (importante hub farmaceutico e biomedicale). Tale polarizzazione riflette congiuntamente la densità degli stabilimenti produttivi regolamentati, la presenza di centri di importazione e distribuzione e la specializzazione industriale regionale. \todo{Inserire i valori esatti dei primi 5 Stati per numero di eventi di richiamo dalla vista cartografica di Dashboard 1.}

\subsection{Profilo di Letalità per Settore Merceologico}
L'analisi incrociata tra \texttt{Tipo Prodotto (Atomico)} ed \texttt{Event Classification} (``Profilo di Letalità per Settore Merceologico'') mette in luce differenze strutturali marcate nella distribuzione percentuale tra Classe~I, Classe~II e Classe~III, sintetizzate nella Tabella~\ref{tab:sector_lethality}.

\begin{table}[htbp]
\centering
\caption{Profilo di letalità per settore merceologico (incidenza percentuale per Classe di Rischio FDA rilevata nella Dashboard~1).}
\label{tab:sector_lethality}
\begin{tabularx}{\textwidth}{@{}Xrrr@{}}
\toprule
\textbf{Settore Merceologico} & \textbf{Classe I (\%)} & \textbf{Classe II (\%)} & \textbf{Classe III (\%)} \\
\midrule
Prodotti per Animali & 47,65 & 39,96 & 12,39 \\
Cibo e Bevande & $\approx$ 43,57 & $\approx$ 50,43 & Quota residua \\
Medicinali & 9,70 & 80,76 & 9,54 \\
Dispositivi & 9,09 & 88,33 & Quota residua \\
Cosmetici & Quota residua & 92,94 & Quota residua \\
Prodotti Biologici & Quota marginale & 67,70 & 31,93 \\
Tabacco & 0,00 & 100,00 & 0,00 \\
\bottomrule
\end{tabularx}
\end{table}

I settori legati alla filiera alimentare e mangimistica mostrano la maggiore incidenza di rischio clinico estremo: nei \textit{Prodotti per Animali} la Classe~I raggiunge il 47,65\% (superando la Classe~II al 39,96\% e la Classe~III al 12,39\%), mentre in \textit{Cibo e Bevande} la Classe~I si attesta attorno al 43,57\% a fronte di circa il 50,43\% di Classe~II. Al contrario, \textit{Dispositivi} (88,33\% Classe~II, 9,09\% Classe~I), \textit{Medicinali} (80,76\% Classe~II, 9,70\% Classe~I, 9,54\% Classe~III) e \textit{Cosmetici} (92,94\% Classe~II) sono dominati dalla Classe~II. I \textit{Prodotti Biologici} presentano una quota particolarmente elevata di richiami procedurali di Classe~III (31,93\%) accanto al 67,70\% di Classe~II, mentre il \textit{Tabacco} ricade integralmente in Classe~II (100,00\%).

\subsection{Classifica Top 15 Aziende per Volume e Gravità}
Il pannello ``Classifica Top 15 Aziende per Volume e Gravità'' ordina i primi 15 operatori per numero di prodotti richiamati (\texttt{Conteggio prodotti}), scomponendo ciascuna barra per classe di severità:
\begin{itemize}
    \item \textbf{Prevalenza di operatori medicali, diagnostici ed ematologici (dominati da Classe~II/III)}: figurano ai vertici \textit{ICU Medical} (837 Classe~II, 87 Classe~I e una quota marginale in Classe~III), \textit{Zimmer Biomet} (872 Classe~II), \textit{Medline Industries LP - Northfield} (750 Classe~II e 72 Classe~I), \textit{Siemens Medical Solutions USA} (759 Classe~II), \textit{OneBlood} (664 Classe~II e 103 Classe~III), \textit{Siemens Healthcare Diagnostics} (643 Classe~II e 43 Classe~III), \textit{Boston Scientific} (507 Classe~II e 141 Classe~I), \textit{American National Red Cross} (455 Classe~II e 189 Classe~III), \textit{GE Healthcare} (522 Classe~II e 42 Classe~I), \textit{Medline Industries LP} (469 Classe~II e 87 Classe~I) e \textit{Aidapak Services} (538 Classe~II).
    \item \textbf{Operatori ad altissima incidenza di Classe~I}: si distinguono per criticità clinica \textit{Customed} (412 Classe~I e 270 Classe~II), \textit{Newly Weds Foods} (432 Classe~II e 220 Classe~I), \textit{Garden-Fresh Foods} (610 Classe~I) e \textit{King Bio} (465 Classe~I).
\end{itemize}
Si rileva inoltre che la presenza separata di \textit{Medline Industries LP - Northfield} e \textit{Medline Industries LP} testimonia l'impatto della granularità delle sedi operative sul calcolo delle graduatorie aziendali, nonostante la normalizzazione introdotta da \texttt{Ragione Sociale (Pulita)}.

\keyinsight{Evidenza Chiave -- Dashboard 1}{La severità clinica non coincide con il mero volume aggregato: mentre i settori dei dispositivi medici e dei farmaci generano i maggiori volumi assoluti presso grandi gruppi industriali (prevalentemente in Classe~II), i comparti \textit{Prodotti per Animali} (47,65\% Classe~I) e \textit{Cibo e Bevande} ($\approx$ 43,57\% Classe~I) concentrano la quota proporzionalmente più elevata di rischi critici per la salute, riflessa nella presenza di aziende alimentari e farmaceutiche ad alta quota di Classe~I tra le prime 15 posizioni.}

\section{Dashboard 2: Analisi Temporale e Tassonomia delle Cause di Richiamo}
\label{subsec:dash2}

La seconda dashboard del progetto DashMe è dedicata alla domanda di ricerca \textbf{RQ2}, con l'obiettivo di esaminare l'andamento diacronico degli eventi regolatori, la scomposizione delle cause radice e l'incidenza specifica degli allergeni non dichiarati. Poiché la vista grafica definitiva non è ancora stata esportata in formato immagine, la Figura~\ref{fig:dash2} funge da segnaposto strutturato nell'attesa dell'inclusione della tavola finale.

\dashboardfig{Dash2.png}{Dashboard 2: Analisi Temporale e Tassonomia delle Cause di Richiamo.}{fig:dash2}

Per garantire piena coerenza analitica con le altre due dashboard, il completamento di questa sezione dovrà articolarsi lungo tre assi di indagine:
\begin{enumerate}
    \item \textbf{Evoluzione temporale e stagionalità}: analisi della serie storica basata su \texttt{Center Classification Date} per identificare eventuali trend pluriennali, picchi anomali associati a grandi richiami multi-prodotto e ciclicità stagionali (es.\ nei comparti alimentari freschi). \todo{Esportare Dash2.png da Tableau nella cartella docs/Tableau/ e descrivere l'andamento temporale (anni di picco, trend medio e stagionalità mensile) dei richiami.}
    \item \textbf{Tassonomia delle cause di richiamo}: quantificazione del peso relativo delle categorie estratte nel campo calcolato \texttt{Categoria Ragione di Richiamo} (contaminazioni microbiologiche da Salmonella, Listeria ed E.Coli, allergeni non dichiarati, corpi estranei, difetti di sterilità o etichettatura, adulterazioni farmacologiche), distinguendo tra frequenza per evento e impatto in termini di \texttt{Conteggio prodotti}. \todo{Inserire la graduatoria quantitativa delle macro-cause in Categoria Ragione di Richiamo e la relativa incidenza per Classe I, II e III.}
    \item \textbf{Focus sugli allergeni}: approfondimento mediante i campi calcolati \texttt{Dettaglio Allergene} e \texttt{Numero Allergeni Coinvolti} per individuare gli allergeni più ricorrenti (es.\ latte, frutta a guscio, soia, grano, uova, arachidi) e valutare la frequenza dei richiami multi-allergene. \todo{Riportare i valori estratti da Dettaglio Allergene e la distribuzione di Numero Allergeni Coinvolti dalla Dashboard 2.}
\end{enumerate}

\section{Dashboard 3: Governance Operativa, Portata Distributiva, Conservazione e Packaging}
\label{subsec:dash3}

La terza dashboard (\textit{``Governance Operativa, Portata Distributiva, Conservazione e Packaging''}, Figura~\ref{fig:dash3}) affronta la domanda \textbf{RQ3}, spostando il fuoco dall'origine geografica e settoriale alla gestione operativa del richiamo lungo la filiera logistica e di confezionamento.

\dashboardfig{Dash3.png}{Dashboard 3: Governance Operativa, Portata Distributiva, Conservazione e Packaging.}{fig:dash3}

\subsection{Quote di Richiami sul Mercato e Stato di Chiusura}
Il pannello ``Quote di Richiami sul Mercato'' analizza la distribuzione di 104.537 prodotti richiamati complessivi rispetto al campo \texttt{Status}:
\begin{itemize}
    \item \textbf{Terminato} (\textit{Terminated}): 84.502 prodotti;
    \item \textbf{In corso} (\textit{Ongoing}): 18.555 prodotti;
    \item \textbf{Completato} (\textit{Completed}): 1.480 prodotti.
\end{itemize}
Sebbene la maggioranza dei prodotti rientri in procedure formalmente terminate, la persistenza di 18.555 prodotti in stato ``In corso'' rappresenta un carico operativo rilevante per il monitoraggio post-vendita e la verifica dell'effettivo ritiro dai canali distributivi.

\subsection{Ruolo della Catena del Freddo}
La vista ``Ruolo della Catena del Freddo'' incrocia il campo calcolato \texttt{Modalità Conservazione} con \texttt{Categoria Ragione di Richiamo}. In termini volumetrici, la categoria \textit{Temperatura Ambiente / Altro} concentra circa 38.000 prodotti richiamati, a fronte di circa 3.000 prodotti nel segmento \textit{Refrigerato / Fresco} e circa 2.500 nel segmento \textit{Surgelato}. Nonostante i volumi inferiori rispetto alla fascia a temperatura ambiente (che ingloba farmaci, dispositivi e prodotti secchi), le barre impilate mostrano come i regimi refrigerato e surgelato siano esposti in modo specifico a rischi biologici e di etichettatura --- in particolare \textit{Listeria}, \textit{Salmonella}, \textit{E.Coli}, allergeni non dichiarati e corpi estranei --- mentre nella fascia a temperatura ambiente emergono anche le adulterazioni dei farmaci accanto alle contaminazioni e agli allergeni. \todo{Verificare e riportare la ripartizione esatta dei volumi per ciascuna causa all'interno delle tre fasce di Modalità Conservazione.}

\subsection{Correlazione tra Packaging e Rischi}
Il grafico ``Correlazione tra Packaging e Rischi'' mette in relazione \texttt{Tipo Packaging} con \texttt{Categoria Ragione di Richiamo}. Tra le tipologie esplicitamente identificate dal parser testuale spiccano:
\begin{itemize}
    \item \textbf{Plastica / Busta}: oltre 8.000 prodotti richiamati;
    \item \textbf{Vetro / Bottiglia / Barattolo}: circa 4.800 prodotti;
    \item \textbf{Cartone / Scatola}: circa 2.300 prodotti;
    \item \textbf{Lattina / Metallo}: circa 2.100 prodotti.
\end{itemize}
Le categorie \textit{Plastica / Busta} e \textit{Vetro / Bottiglia / Barattolo} risultano dominate da richiami per \textit{Salmonella}, allergeni non dichiarati e \textit{Listeria}, segnalando la stretta interdipendenza tra imballaggi primari di prodotti alimentari freschi o pronti al consumo e rischi di contaminazione microbiologica o di errore nel confezionamento/etichettatura.

\subsection{Matrice Portata Distributiva vs Severità Clinica}
Infine, la matrice ``Portata Distributiva vs Severità Clinica'' incrocia i quattro livelli di \texttt{Portata Distributiva} (\textit{Internazionale / Export}, \textit{Nazionale (Tutti gli USA)}, \textit{Regionale (Multi-Stato)}, \textit{Locale (Singolo Stato)}) con le tre classi FDA (\textit{Classe I}, \textit{Classe II}, \textit{Classe III}). Complessivamente la \textit{Classe~II} costituisce il segmento numericamente predominante attraverso le diverse ampiezze di distribuzione; tuttavia, la matrice evidenzia una presenza significativa di richiami in \textit{Classe~I} proprio nelle fasce a maggiore estensione geografica e complessità logistica, segnatamente \textit{Internazionale / Export} e \textit{Regionale (Multi-Stato)}. \todo{Inserire i conteggi esatti delle 12 celle della matrice Portata Distributiva vs Severità Clinica direttamente dai tooltip di Tableau.}

\keyinsight{Evidenza Chiave -- Dashboard 3}{La vulnerabilità operativa dei richiami non risiede soltanto nella produzione, ma nell'intersezione tra logistica, imballaggio e governance di chiusura: 18.555 prodotti risultano ancora ``In corso'' su 104.537 totali, i rischi microbiologici e allergenici si concentrano su imballaggi in plastica/busta ($>8\text{K}$) e vetro ($\approx 4,8\text{K}$), e una quota significativa di eventi critici di Classe~I interessa catene distributive regionali multi-stato e internazionali.}

\section{Discussione: Sintesi Trasversale, Implicazioni Manageriali e Limiti}
\label{sec:discussion}

\subsection{Sintesi Trasversale dei Risultati}
La lettura integrata delle viste sviluppate in DashMe consente di superare l'analisi per singoli indicatori isolati, evidenziando tre dinamiche strutturali tra loro interconnesse. In primo luogo, emerge una chiara dissociazione tra \textit{volume assoluto} e \textit{criticità clinica}: i comparti dei dispositivi medici e dei prodotti farmaceutici guidano i volumi complessivi di richiamo presso grandi operatori industriali, ma rimangono ancorati prevalentemente alla Classe~II (rispettivamente 88,33\% e 80,76\%); di converso, i comparti \textit{Prodotti per Animali} (47,65\% Classe~I) e \textit{Cibo e Bevande} ($\approx$ 43,57\% Classe~I) concentrano la quota proporzionalmente più severa di eventi potenzialmente letali.

In secondo luogo, la dimensione geografica e quella della portata distributiva mostrano come il rischio non resti confinato all'area di origine. Se gli Stati di California, Texas, Florida, il Nord-Est/Midwest e Puerto Rico rappresentano i principali nodi di notifica (fino a 4.522 eventi per singolo Stato), la matrice di Dashboard~3 conferma che una quota rilevante di richiami di Classe~I raggiunge circuiti \textit{Regionali (Multi-Stato)} e \textit{Internazionali / Export}, amplificando la complessità del coordinamento inter-giurisdizionale.

In terzo luogo, l'osservazione congiunta delle cause di richiamo, della conservazione e del confezionamento suggerisce che contaminazioni microbiologiche (\textit{Salmonella}, \textit{Listeria}, \textit{E.Coli}) e allergeni non dichiarati trovano un vettore ricorrente negli imballaggi in \textit{Plastica / Busta} (oltre 8.000 prodotti) e \textit{Vetro / Bottiglia / Barattolo} (circa 4.800 prodotti), interessando sia il segmento a temperatura ambiente sia le filiere \textit{Refrigerato / Fresco} e \textit{Surgelato}. \todo{Integrare la sintesi trasversale con le dinamiche temporali e la classifica degli allergeni non appena completata l'analisi della Dashboard 2.}

\subsection{Implicazioni per il Project Management e il Risk Management}
Dal punto di vista gestionale e dell'ingegneria dei processi di progetto, i risultati offrono indicazioni operative concrete per i responsabili della qualità, della \textit{supply chain} e della conformità regolatoria:
\begin{enumerate}
    \item \textbf{Audit selettivi sui fornitori e allocazione delle risorse di controllo}: la matrice di rischio emersa dalla Dashboard~1 suggerisce di differenziare la frequenza e la profondità degli audit in base al profilo di letalità settoriale e alla concentrazione territoriale. Per fornitori attivi in \textit{Cibo e Bevande}, \textit{Prodotti per Animali} o presso siti ad alta incidenza di Classe~I, gli audit dovrebbero privilegiare test microbiologici e controlli sui cambi formato/etichetta; nei dispositivi medici e nei centri ematologici, invece, l'attenzione va posta sul controllo statistico di processo e sulla gestione documentale delle non conformità ricorrenti.
    \item \textbf{Governance della logistica inversa e chiusura dei 18.555 richiami ``In corso''}: la presenza di 18.555 prodotti in stato ``In corso'' su 104.537 totali rilevati nella Dashboard~3 evidenzia come l'apertura formale di un richiamo non garantisca una rapida chiusura operativa. Nei progetti di gestione delle crisi (\textit{recall management}), è indispensabile trattare ogni evento come un programma strutturato con \textit{milestone} di tracciabilità, riconciliazione delle scorte rese o distrutte e verifica della ricezione degli avvisi presso i distributori, soprattutto quando la portata è multi-stato o internazionale.
    \item \textbf{Controlli integrati su catena del freddo e linee di confezionamento}: l'esposizione di \textit{Plastica / Busta} e \textit{Vetro / Bottiglia / Barattolo} a richiami per \textit{Salmonella}, \textit{Listeria} e allergeni non dichiarati richiama l'esigenza di rafforzare i punti critici di controllo (HACCP) proprio nelle fasi finali di riempimento, sigillatura ed etichettatura, oltre al monitoraggio continuo delle escursioni termiche nei segmenti \textit{Refrigerato / Fresco} e \textit{Surgelato}.
\end{enumerate}

\subsection{Limiti Metodologici}
Una valutazione rigorosa dei risultati richiede di esplicitare quattro limiti intrinseci al dataset e alla pipeline analitica adottata:
\begin{itemize}
    \item \textbf{Tassonomia basata su regole testuali e parole chiave}: i campi calcolati \texttt{Categoria Ragione di Richiamo}, \texttt{Tipo Packaging}, \texttt{Modalità Conservazione}, \texttt{Portata Distributiva} e \texttt{Dettaglio Allergene} operano tramite corrispondenze lessicali su stringhe in linguaggio naturale. Pur essendo efficaci per l'aggregazione esplorativa, tali regole possono risentire di sinonimi non mappati, descrizioni ambigue o compresenza di più concetti nello stesso testo.
    \item \textbf{Ampiezza della classe residuale ``Altro / Non specificato''}: nei campi derivati da testo libero (in particolare nella conservazione a temperatura ambiente/altro e nel packaging), una quota significativa di record non contiene indicazioni esplicite nel testo sorgente e confluisce nella categoria residuale. \todo{Quantificare esattamente l'incidenza percentuale della classe residuale ``Altro / Non specificato'' per Tipo Packaging, Modalità Conservazione e Categoria Ragione di Richiamo.}
    \item \textbf{Natura amministrativa e auto-dichiarata dei dati FDA}: il registro riflette tempistiche di classificazione (\texttt{Center Classification Date}) e aggiornamenti di stato (\texttt{Status}) che possono subire ritardi amministrativi rispetto al momento effettivo del ritiro fisico dal mercato. Inoltre, l'unità di conteggio adottata (\texttt{Conteggio prodotti}) pondera maggiormente eventi multi-referenza rispetto a richiami mono-prodotto.
    \item \textbf{Residua frammentazione anagrafica}: nonostante il campo \texttt{Ragione Sociale (Pulita)}, operatori multi-sito come \textit{Medline Industries LP} e \textit{Medline Industries LP - Northfield} possono ancora apparire come entità distinte nella classifica Top~15 quando i record mantengono riferimenti di stabilimento differenti.
\end{itemize}

\section{Conclusioni e Sviluppi Futuri}
\label{sec:conclusions}

Il progetto \textbf{DashMe} dimostra come l'applicazione sistematica di tecniche di \textit{Business Intelligence}, modellazione ETL in Tableau e \textit{Visual Analytics} permetta di trasformare un archivio regolatorio complesso e testualmente eterogeneo come quello dei richiami FDA in un sistema decisionale leggibile e azionabile. Partendo dalle 17 colonne grezze di \texttt{data/recalls.xlsx} e costruendo 14 campi calcolati in \texttt{src/Tableau/dashsmash3.twb} su estrazione \texttt{data/dati\_trasformati\_tableau.hyper}, il lavoro consente di rispondere alle tre domande di ricerca formulate:

\begin{itemize}
    \item \textbf{Risposta a RQ1 (Geografia, settori e responsabilità aziendale)}: gli eventi di richiamo risultano fortemente concentrati in specifici poli produttivi e logistici statunitensi (fino a 4.522 eventi per singolo Stato, con preminenza di California, Texas, Florida, Nord-Est/Midwest e Puerto Rico). A livello settoriale, i comparti medicali, farmaceutici e cosmetici presentano una prevalenza di Classe~II (80,76\%--92,94\%, fino al 100,00\% nel tabacco), mentre \textit{Prodotti per Animali} (47,65\% Classe~I) e \textit{Cibo e Bevande} ($\approx$ 43,57\% Classe~I) concentrano la massima severità clinica. Nella graduatoria Top~15 per volume emergono grandi operatori medicali ed ematologici accanto ad aziende alimentari e farmaceutiche con elevatissima incidenza di Classe~I.
    \item \textbf{Risposta a RQ2 (Dinamiche temporali e tassonomia delle cause)}: l'architettura semantica per l'analisi temporale e causale è stata predisposta mediante i campi \texttt{Categoria Ragione di Richiamo}, \texttt{Dettaglio Allergene} e \texttt{Numero Allergeni Coinvolti}; la consolidazione empirica di trend, stagionalità e profili allergenici è formalmente strutturata nella Sezione~\ref{subsec:dash2} in attesa dell'esportazione definitiva di \texttt{Dash2.png}. \todo{Aggiornare la sintesi conclusiva di RQ2 con i risultati quantitativi estratti da Dashboard 2.}
    \item \textbf{Risposta a RQ3 (Governance operativa, catena del freddo, packaging e distribuzione)}: su 104.537 prodotti richiamati, 84.502 risultano terminati e 1.480 completati, ma ben 18.555 permangono ``In corso'', segnalando un'area critica di gestione post-notifica. Sul piano operativo, gli imballaggi in \textit{Plastica / Busta} ($>$8.000) e \textit{Vetro / Bottiglia / Barattolo} ($\approx$ 4.800) concentrano i richiami da \textit{Salmonella}, allergeni e \textit{Listeria}, mentre la matrice distributiva evidenzia come una quota significativa di eventi di Classe~I interessi reti di vendita regionali multi-stato e internazionali.
\end{itemize}

Il contributo principale di DashMe risiede nell'aver colmato il divario tra dato regolatorio grezzo e lettura orientata al \textit{Project Management} e al \textit{Risk Management}, offrendo una struttura analitica riproducibile per la prioritizzazione di audit, controlli di packaging e piani di ritiro.

In prospettiva futura, il progetto si presta a tre direttrici di evoluzione:
\begin{enumerate}
    \item \textbf{Finalizzazione della Dashboard~2}: completare l'esportazione grafica e la validazione quantitativa delle serie storiche, della stagionalità e della scomposizione multi-allergene;
    \item \textbf{Integrazione dinamica tramite API openFDA}: sostituire lo snapshot statico su foglio elettronico con una pipeline automatizzata verso gli endpoint ufficiali openFDA, abilitando l'aggiornamento incrementale dell'estrazione \texttt{.hyper} e il monitoraggio quasi in tempo reale dei richiami ``In corso'';
    \item \textbf{Modellazione predittiva della severità}: utilizzare le feature ingegnerizzate (settore merceologico, macro-causa, packaging, modalità di conservazione e portata distributiva) per addestrare modelli di classificazione capaci di stimare precocemente la probabilità di ricaduta in Classe~I al momento della segnalazione iniziale.
\end{enumerate}


\bibliographystyle{unsrtnat}
\bibliography{references}

\end{document}
